\documentclass[a4paper]{article}

\usepackage[utf8x]{inputenc}    
\usepackage[T1]{fontenc}
\usepackage{graphicx}
\usepackage[frenchb]{babel}
\usepackage[francais,bloc,nowatermark]{automultiplechoice}    
% [nowatermark] pour enlever le texte "PROJET" et le pied de page
\usepackage{multicol}
\usepackage{tkz-tab}
\usepackage{fancyhdr,amssymb,amsmath}
%\DeclareGraphicsRule{.7}{mps}{*}{}

%===================
\begin{document}
%===================

%%% On répartie les quesitons dans différents groupes.
% ---------------------------------
%%% Questions groupe : "catégorie1"
%----------------------------------
\element{categorie1}{
  \begin{question}{Derive1}    
% Question ouverte   
Soit $f$ une fonction définie sur $\mathbb{R}$ telle que $f(x)=5x^3-4x^2+x-3$. Déterminer $f'(x)$, puis calculer $f(1)$ et $f'(1)$.
   
   \AMCOpen{lines=3,dots=false}{\wrongchoice{NA}\scoring{0}\wrongchoice{AB}\scoring{1}\wrongchoice{B}\scoring{2}\correctchoice{TB}\scoring{3}}
% AB rapporte 1 point
% B rapporte 2 points
% TB rapporte 3 points     
  \end{question}
}

%
\element{categorie1}{
  \begin{question}{Derive2}  
% Question ouverte  
Soit $f$ une fonction définie sur $\mathbb{R}$ telle que $f(x)=-3x^3+4x^2-x-3$. Déterminer $f'(x)$, puis calculer $f(2)$ et $f'(2)$.
   \AMCOpen{lines=3,dots=false}{\wrongchoice{NA}\scoring{0}\wrongchoice{AB}\scoring{1}\wrongchoice{B}\scoring{2}\correctchoice{TB}\scoring{3}}
% AB rapporte 1 point
% B rapporte 2 points
% TB rapporte 3 points         
  \end{question}
}

\element{categorie1}{
  \begin{question}{Derive3}  
% Question ouverte  
Soit $f$ une fonction définie sur $\mathbb{R}$ telle que $f(x)=-3x^3+4x^2-5x-3$. Déterminer $f'(x)$, puis calculer $f(-1)$ et $f'(-1)$.
   \AMCOpen{lines=3,dots=false}{\wrongchoice{NA}\scoring{0}\wrongchoice{AB}\scoring{1}\wrongchoice{B}\scoring{2}\correctchoice{TB}\scoring{3}}
% AB rapporte 1 point
% B rapporte 2 points
% TB rapporte 3 points         
  \end{question}
}

\element{categorie2}{
  \begin{question}{variation1}  
% Question ouverte  
Soit $f$ une fonction définie sur $\mathbb{R}$ telle que $f(x)=x^3+5x-10$. Déterminer les variations de $f$ sur $\mathbb{R}$ à l'aide de la fonction dérivée et d'un tableau récapitulatif.
   \AMCOpen{lines=8,dots=false}{\wrongchoice{NA}\scoring{0}\wrongchoice{AB}\scoring{1}\wrongchoice{B}\scoring{2}\correctchoice{TB}\scoring{3}}
% AB rapporte 1 point
% B rapporte 2 points
% TB rapporte 3 points         
  \end{question}
}

\element{categorie2}{
  \begin{question}{variation2}  
% Question ouverte  
Soit $f$ une fonction définie sur $\mathbb{R}$ telle que $f(x)=-2x^3-4x+5$. Déterminer les variations de $f$ sur $\mathbb{R}$ à l'aide de la fonction dérivé et d'un tableau récapitulatif.
   \AMCOpen{lines=8,dots=false}{\wrongchoice{NA}\scoring{0}\wrongchoice{AB}\scoring{1}\wrongchoice{B}\scoring{2}\correctchoice{TB}\scoring{3}}
% AB rapporte 1 point
% B rapporte 2 points
% TB rapporte 3 points         
  \end{question}
}

\element{categorie3}{
  \begin{question}{tangente1}  
% Question ouverte  
On rappelle que l'équation réduite de la tangente au point d'abscisse $a$ d'une fonction $f$ est donnée par $\Gamma_{a,f}: y=f'(a)(x-a)+f(a)$ 
Soit $h$ une fonction définie sur $\mathbb{R}$ telle que $h(x)=x^2-x+5$. Déterminer l'équation réduite de $\Gamma_{-2,h}$. 
   \AMCOpen{lines=3,dots=false}{\wrongchoice{NA}\scoring{0}\wrongchoice{AB}\scoring{1}\wrongchoice{B}\scoring{2}\correctchoice{TB}\scoring{3}}
% AB rapporte 1 point
% B rapporte 2 points
% TB rapporte 3 points         
  \end{question}
}


\element{categorie3}{
  \begin{question}{tangente2}  
% Question ouverte  
On rappelle que l'équation réduite de la tangente au point d'abscisse $a$ d'une fonction $f$ est donnée par $\Gamma_{a,f}: y=f'(a)(x-a)+f(a)$ 
Soit $h$ une fonction définie sur $\mathbb{R}$ telle que $h(x)=x^3-x+5$. Déterminer l'équation réduite de $\Gamma_{-1,h}$. 
   \AMCOpen{lines=3,dots=false}{\wrongchoice{NA}\scoring{0}\wrongchoice{AB}\scoring{1}\wrongchoice{B}\scoring{2}\correctchoice{TB}\scoring{3}}
% AB rapporte 1 point
% B rapporte 2 points
% TB rapporte 3 points         
  \end{question}
}



\element{categorie3}{
  \begin{question}{tangente3}  
% Question ouverte  

\textit{(On rappelle que l'équation réduite de la tangente au point d'abscisse $a$ d'une fonction $f$ est donnée par $\Gamma_{a,f}: y=f'(a)(x-a)+f(a)$)} \\
Soit $h$ une fonction définie sur $\mathbb{R}$ telle que $h(x)=x^3-2x-6$. Déterminer l'équation réduite de $\Gamma_{-2,h}$. 
   \AMCOpen{lines=3,dots=false}{\wrongchoice{NA}\scoring{0}\wrongchoice{AB}\scoring{1}\wrongchoice{B}\scoring{2}\correctchoice{TB}\scoring{3}}
% AB rapporte 1 point
% B rapporte 2 points
% TB rapporte 3 points         
  \end{question}
}

\element{categorie4}{
  \begin{question}{libre1}  
% Question ouverte  
Représenter graphiquement une fonction $f$ de votre choix sur l'intervalle $[-3;5]$, tracer la tangente au point d'abscisse $2$ de la courbe $\mathcal{C}_f$ et indiquer le signe ou la valeur de $f'(2)$.
   \AMCOpen{lines=12,dots=false}{\wrongchoice{NA}\scoring{0}\wrongchoice{AB}\scoring{1}\wrongchoice{B}\scoring{2}\correctchoice{TB}\scoring{3}}
% AB rapporte 1 point
% B rapporte 2 points
% TB rapporte 3 points         
  \end{question}
}


\element{categorie4}{
  \begin{question}{libre2}  
% Question ouverte  
Soit $i$ une fonction dont le tableau de signe de $i'$ à été représenté sur $[-3;10]$. 
\begin{center}
\begin{tikzpicture}
   \tkzTabInit{$x$ / 1 , $f'(x)$ / 1}{$-3$, $6$, $8$, $10$}
   \tkzTabLine{, -, z, +, z, -, }
\end{tikzpicture}
\end{center}
Représenter graphiquement une fonction qui pourrait correspondre à $i$ sur l'intervalle $[-3;10]$.


   \AMCOpen{lines=12,dots=false}{\wrongchoice{NA}\scoring{0}\wrongchoice{AB}\scoring{1}\wrongchoice{B}\scoring{2}\correctchoice{TB}\scoring{3}}
% AB rapporte 1 point
% B rapporte 2 points
% TB rapporte 3 points         
  \end{question}
}

%-------------------------

%%%%%%%%%%%%%%%%%%%%%%%%%
%% fin des questions %%%%
%%%%%%%%%%%%%%%%%%%%%%%%%


%%%%%%%%%%%%%%%%%%%%%%%%%%%%
%%% fabrication des copies%%
%%%%%%%%%%%%%%%%%%%%%%%%%%%%
\exemplaire{20}{    

%%% debut de l'entête des copies :    
	%%% première ligne

	\begin{minipage}{1\linewidth}
	  \centering\large\bf DS de Mathématiques 1STMG sur les fonctions dérivées    
	\end{minipage}

	%%% codage des numéros d'étudiants sur 3 chiffres par les candidats et encadré d'écriture manuelle des noms.
%	\vspace*{3mm}
	{
	%\setlength{\parindent}{0pt}\AMCcode{etu}{2}\hspace*{\fill}
	\begin{minipage}[b]{12cm}
	% \hspace*{5mm}
  	%  Codez votre numéro ci contre chiffre par chiffre, puis complétez l'encadré.
    % \vspace{3ex}

	 \hfill\champnom{
	 \centering
	 \fbox{
	 \begin{minipage}{1\linewidth}
		\vspace*{3mm}
		NOM - Prénom - Classe :
		\vspace*{3mm}
	 \end{minipage}
	    }	
	   }
	\hfill\vspace{1ex}
% texte de présentation 
\begin{center}\em

Durée : 55 minutes.

  Aucun document n'est autorisé. La calculatrice est autorisée.
%  Les questions faisant apparaître le symbole \multiSymbole{} peuvent
%  présenter une ou plusieurs bonnes réponses. Les autres ont une unique bonne réponse.
%  Les réponses fausses ou incohérentes retirent des points. Pour éviter les erreurs machines, il est nécessaire de \textbf{noircir vos cases}.

\end{center}

	\end{minipage}\hspace*{\fill}
	}

\vspace{4ex}

%%% Fin de l'entête
%%%%%%%%%%%%%%%%%%%%%%%

%% Composition des QCM
%%% On mélange des questions par groupe
% On efface le {tout} précédent
\cleargroup{tout}

% On mélange par "catégorie", et on prélève aléatoirement [n] questions 
% du lot {categorie-k} vers la copie complète {tout}
% puis on restitue la copie.
% Si on met [0], toutes les questions sont prises.
\melangegroupe{categorie1}\copygroup[3]{categorie1}{tout}
\melangegroupe{categorie2}\copygroup[1]{categorie2}{tout}
\melangegroupe{categorie3}\copygroup[1]{categorie3}{tout}
\melangegroupe{categorie4}\copygroup[0]{categorie4}{tout}

\restituegroupe{tout}

% Saut d'une page si le sujet comporte un nombre impair de pages.
% Ainsi il est possible d'imprimer les sujets en recto-verso.

\AMCcleardoublepage    

%%% Fin de la création des questionnaires

}
% Fin de exemplaires{k}{ 

\end{document}

